% 辅助线方法 C：延长两腰相交于 P，得 A 字相似
\documentclass[tikz,border=4pt]{standalone}
\usepackage{ctex}
\usepackage{amsmath}
\usetikzlibrary{calc}
\begin{document}
\begin{tikzpicture}[scale=1.0, thick]
  % 选定让两腰延长后在上方交于 P
  \coordinate (B) at (0,0);
  \coordinate (C) at (5,0);
  \coordinate (A) at (1.5,2);
  \coordinate (D) at (3.5,2);
  % 直线 BA：从 B 沿方向 (1.5,2) 延长；直线 CD：从 C 沿方向 (-1.5,2) 延长
  % 交点 P：B+t(1.5,2)=C+s(-1.5,2)，解得 t=s=5/3，P=(2.5, 10/3)
  \coordinate (P) at (2.5,3.333);

  \draw (A)--(B)--(C)--(D)--cycle;
  \draw[red, dashed] (A)--(P);
  \draw[red, dashed] (D)--(P);

  \node[left]  at (A) {$A$};
  \node[right] at (D) {$D$};
  \node[below left]  at (B) {$B$};
  \node[below right] at (C) {$C$};
  \node[above] at (P) {$P$};

  \node[below] at (2.5,-0.6) {\footnotesize 方法 C：延长两腰};
\end{tikzpicture}
\end{document}
